\documentclass{situationist-preprint}

\title{The Dispute Is Notational}
\msid{ETH-02-c}
\msfamily{ETH-02}
\msdiscipline{Normative and Applied Ethics}
\mstype{Consequence}
\mskeywords{consequentialism; deontology; consequentializing; notational variance; agent-relative value; virtue ethics; methodology}

\begin{document}
\maketitle

\begin{abstract}
ETH-02-a proves that deontological and consequentialist theories are extensionally equivalent, with mutually inverse translations. We draw the consequences for the practice of normative ethics. No verdict about a case, or any finite set of verdicts, can favor one family over the other; arguments of the form \enquote{this family gives the wrong verdict here} refute at most a member of the family. The disagreements that survive are first-order disagreements about verdicts, which cut across the families, and one metaethical commitment, the intelligibility of agent-relative value, on which the notational verdict depends. We show that the same reasoning extends to virtue ethics as characterized in ETH-12.
\end{abstract}

\section{What the equivalence rules out}

ETH-02-a proves that, given act-inclusive outcomes, a deontic theory is consequentializable if and only if it is admissible and has no dilemmas (Theorem~3.5), that every consequentialist theory is such a theory (Theorem~3.4), and that the translations between the two forms of presentation are mutually inverse up to extensional equivalence (Corollary~3.6). Its Section~6 adopts a criterion on which two theories are notational variants with respect to a subject matter if each can be translated into the other, preserving all its claims about that subject matter, without further substantive input. We draw three consequences.

\begin{proposition}\label{p:cases}
Let $K$ be any finite set of verdicts, each assigning permissibility or impermissibility to an option for an agent in a situation, such that the verdicts permit at least one option in each situation they cover. If some deontological theory delivers every verdict in $K$, so does some consequentialist theory, and conversely.
\end{proposition}

\begin{proof}
A deontological theory that delivers $K$ is admissible and without dilemmas (Definition~2.8 of ETH-02-a), and by Theorem~3.3 of ETH-02-a it is equivalent to the consequentialist theory of its indicator value function, which therefore delivers $K$. Conversely, a consequentialist theory that delivers $K$ is, by Theorem~3.4 of ETH-02-a, a deontological theory with a single rule.
\end{proof}

Proposition~\ref{p:cases} has a direct methodological consequence. The traditional argument against consequentialism runs: consequentialism permits pushing the man in Footbridge, or punishing the innocent to prevent a riot, or breaking a promise to secure slightly more good; these verdicts are false; so consequentialism is false. The argument refutes the theories that deliver those verdicts, typically agent-neutral consequentialism on coarse outcomes. It does not refute consequentialism, since by Proposition~\ref{p:cases} the correct verdicts are delivered by some consequentialist theory. The same holds in reverse for arguments against deontology from its alleged rigidity. Cases discriminate between theories, never between the families.

\begin{corollary}\label{c:cut}
Every first-order disagreement between a deontologist and a consequentialist is also a disagreement that two deontologists, or two consequentialists, could have.
\end{corollary}

\begin{proof}
A first-order disagreement is a disagreement about some verdict. If the deontologist's theory delivers it, so does a consequentialist theory by Proposition~\ref{p:cases}, and that theory and the consequentialist's original theory disagree about it; likewise in the other direction.
\end{proof}

Whether one may divert the trolley but not push the man (ETH-03), whether one may lie to the murderer (ETH-04), how much one must give (ETH-06): each of these questions has a deontological and a consequentialist formulation, and the answer is the same whichever formulation is used. A moral philosopher who wants to know what is right gains nothing from first settling which family to work in.

\section{What survives}

The equivalence leaves three things in place, none of which is a dispute between the families about the right.

\paragraph{Explanatory order.} The consequentialist holds that acts are right because their outcomes are best; the deontologist that outcomes are ranked, if at all, because of what it is right to do. Corollary~3.6 of ETH-02-a shows that the two orders deliver the same verdicts, so no verdict can confirm either. Whether one order is metaphysically privileged is a question about grounding in metaethics, not a question about which acts are right, and the equivalence is compatible with either answer.

\paragraph{Agent-relative value.} The equivalence requires agent-relative value functions (Proposition~4.4 of ETH-02-a). \citet{Schroeder2007} has argued that \enquote{good relative to an agent} is not a notion we understand independently of the verdicts it is meant to explain. If he were right, the consequentialist family would shrink to its agent-neutral members, and the dispute would become substantive. The notational verdict therefore depends on a metaethical commitment: that value for an agent is intelligible, for instance as what the agent has reason to prefer, as MTE-11, which takes reasons to be primitive and other normative notions to be defined in terms of them, allows. This is an instance of the thesis of MTE-14 that metaethical commitments constrain first-order theory. It does not reopen the dispute between the families, since on the view of value adopted it is settled, and on the alternative the dispute is not between consequentialism and deontology but between agent-neutral consequentialism and everything else.

\paragraph{Risk and uncertainty.} Under certainty, the cardinal structure of a value function is idle. Under risk, a consequentialist theory that maximizes expected value needs a cardinal scale, and so does any comparison of theories under moral uncertainty (MTE-16). These choices arise within each family once its verdicts under risk are fixed, and the equivalence of ETH-02-a applies to any finite choice structure whatever the options are, including options evaluated as prospects. So risk introduces new questions of scale, but no new question that divides the families.

\section{A third family}

ETH-12 holds that the right action is what the virtuous agent would characteristically do, and that the virtuous agent can be characterized independently of right action. Whatever the details, this gives a deontic theory: an option is permissible for an agent in a situation if and only if it is something a virtuous agent would characteristically do there. Such a theory permits only available options, and, if a virtuous agent always has something she would characteristically do, it has no dilemmas. By Theorem~3.5 of ETH-02-a, it is consequentializable. In the other direction, a consequentialist theory has a virtue-theoretic statement just in case the virtuous agent, characterized independently, would characteristically choose its maximal options; whether this holds for a given value function is a first-order question that the equivalence does not settle. The first direction establishes that virtue ethics, considered as a theory of right action, is not a third option beyond the dispute; it is a further notation for the same verdicts, distinguished, like the others, by what it takes to explain them.

\begin{objection}
If the dispute were notational, the great figures of normative ethics would have been confused about what they disagreed about. That is implausible.
\end{objection}

\begin{reply}
They disagreed about verdicts, and those disagreements were real; Corollary~\ref{c:cut} shows that they survive the equivalence. What they took to follow from the verdicts, that a whole family of theories was refuted, did not follow. Ross showed that a value function ranking outcomes by impersonal good alone cannot deliver the duty of fidelity \citep{Ross1930}; he did not show, and the equivalence shows that no one could show, that no value function can deliver it. The confusion, such as it was, lay in identifying consequentialism with one of its members.
\end{reply}

\section{Conclusion}

From the consequentializing equivalence it follows that no case decides between consequentialism and deontology, that every first-order disagreement between them is a disagreement within each of them, and that the residual differences concern explanatory order, the metaethics of agent-relative value, and the choice of scales under risk. As a dispute about the right, the dispute is notational.

\begin{thebibliography}{9}\small
\bibitem[Ross(1930)]{Ross1930} Ross, W.~D. (1930). \emph{The Right and the Good}. Clarendon Press.
\bibitem[Schroeder(2007)]{Schroeder2007} Schroeder, M. (2007). Teleology, agent-relative value, and `good'. \emph{Ethics}, 117(2), 265--295.
\end{thebibliography}

\end{document}
